\uuid{S2k3}
\titre{Logique}
\auteur{bodin ; croizat ; sacre}
\organisation{exo7}
\licence{CC BY-NC-SA 4.0}
\datecreate{2022-03-08}
\niveau{L1}
\module{Algèbre}
\chapitre{Logique, ensemble, raisonnement}
\sousChapitre{Logique}
\difficulte{3}
\theme{Logique}
\exo7qcmid{hilisit.logique.011}

\contenu{
\begin{question}
\qclassification{100.01}
Pour quelles phrases, l'assertion est vraie si on remplace "$??$" par "$\exists$", mais est fausse si on remplace "$??$ " par "$\forall$" ?
\begin{answers}
	\good{$??\; n \in \Nn^* \quad n \text{ est pair}$}
	\bad{$??\; n \in \Nn^* \quad n(n+1) \text{ est pair}$}
	\good{$??\; n \in \Nn^* \quad n \text{ et } n + 2 \text{ sont des nombres premiers}$}
	\good{$??\; n \in \Nn^* \quad$ si $n$ n'est pas premier alors $n$ admet au moins deux facteurs premiers distincts}
\end{answers}
\begin{explanations}
Il existe des entiers pairs (par exemple $n=4$), mais tous les entiers ne sont pas pairs (par exemple $n=5$).\\
Pour tous les entiers, $n(n+1)$ est pair.\\
Il existe un entier (par exemple $n=11$) tel que $n$ et $n+2$ soient deux nombres premiers ; mais ce n'est pas vrai pour tous les entiers (par exemple pour $n=13$, $n+2$ n'est pas premier). \\
Il existe un entier (par exemple $n=21$) qui a deux facteurs premiers distincts (ici $3$ et $7$) ; mais ce n'est pas vrai pour tous les entiers (par exemple $n=9$, a pour seul facteur premier $3$).
\end{explanations}
\end{question}
}
